\clearpage

\section{Decision Accuracy - Sensitivity Analysis}
\label{app:da-sensitivity}

\paragraph{Sensitivity to near-tied reference rankings.}
Our definition of decision accuracy (DA) requires the ordering of model variants at a proxy point to exactly match their ordering at the reference scale. This strict criterion may overstate decision failures when two variants have nearly identical performance at the reference point: a small numerical reversal is counted identically to a reversal between substantially different variants. We therefore perform a sensitivity analysis in which reference-scale differences within a tolerance $\epsilon$ are treated as effective ties. Specifically, when the proxy and reference rankings disagree but the absolute performance gap between the two variants at the reference point is at most $\epsilon$, we count the comparison as ranking-preserving. We sweep $\epsilon \in \{0, 0.1, 0.25, 0.5, 1, 2\}$ percentage points, with $\epsilon=0$ corresponding to the original DA definition. The set and denominator of pairwise comparisons remain fixed across thresholds, such that changes in DA arise only from reclassifying near-tied ranking reversals.

\begin{table}[ht]
\centering

\label{tab:da_tolerance}
\begin{tabular}{rccc}
\toprule
$\epsilon$ (pp) &
\shortstack{Pairs within\\ tolerance\\at reference} &
\shortstack{Ranking \\ reversals treated \\ as ties} &
\shortstack{Already\\correct} \\
\midrule
0.1  & 357   & 184   & 173   \\
0.25 & 982   & 436   & 546   \\
0.5  & 1,949 & 755   & 1,194 \\
1.0  & 3,268 & 1,099 & 2,169 \\
2.0  & 4,301 & 1,308 & 2,993 \\
\bottomrule
\end{tabular}
\caption{
\textbf{Sensitivity of decision accuracy to the reference-gap tolerance $\epsilon$.}
``Within tolerance'' counts pairs whose absolute performance difference at the reference scale is at most $\epsilon$.
Only reversed pairs are newly reclassified as correct; pairs whose rankings already agree remain unchanged.
}
\end{table}

\begin{figure}[ht]
    \centering
    \includegraphics[width=\linewidth]{figures/da_sensitivity.png}
    \caption{
    \textbf{Sensitivity of decision accuracy to near-tied reference rankings.}
    Left: DA-size across proxy model sizes. Right: DA-checkpoint across training progress for the 350M model.
    Increasing the tolerance $\epsilon$ raises absolute DA by treating small reference-scale ranking reversals as ties, while largely preserving the shape of the curves.
    For DA-size, the curves remain nearly parallel across tolerances, indicating that accounting for near-ties does not induce stronger scale convergence.
    The 1.7B point is the reference scale and therefore reaches perfect agreement by construction.
    }
    \label{fig:da_tolerance}
\end{figure}

Figure~\ref{fig:da_tolerance} shows that increasing the tolerance $\epsilon$ raises absolute DA, indicating that some apparent ranking failures are near-ties at the reference scale. At $\epsilon=0.5$ pp, 755 reversed comparisons are reclassified as correct, increasing to 1,099 at $1$ pp and 1,308 at $2$ pp (Table~\ref{tab:da_tolerance}). However, the DA-size curves remain nearly parallel across thresholds, and the number of reclassified comparisons is similar across proxy sizes. Thus, tolerance mainly shifts DA upward rather than making larger proxy models more predictive of the reference ranking. DA-checkpoint shows the same upward shift while preserving its increase with training progress. Overall, near-ties affect the absolute DA level but do not change the main conclusion: ranking recovery improves with training progress, but not consistently with model size.

\paragraph{Categorizing changes in intervention effects.}
Decision accuracy measures whether the ordering of model variants is preserved from a proxy point to the reference point, but it does not distinguish between qualitatively different ways in which pairwise differences can evolve. In particular, two variants may remain similarly close at both proxy and reference, an apparent proxy-scale difference may disappear at the reference, or a difference may emerge only at larger scale or later in training. To characterize these cases directly, we compare the absolute performance gap between each pair of variants at the proxy and reference points. We use a practical tolerance of $\epsilon=0.5$ percentage points and treat pairs with an absolute gap of at most $\epsilon$ as ties. This yields four categories: \emph{tie--tie}, where the variants remain within the tolerance at both points; \emph{difference--tie}, where a proxy-scale difference washes out at the reference; \emph{tie--difference}, where a difference emerges only at the reference; and \emph{difference--difference}, where the variants remain separated by more than $\epsilon$ at both points. We compute these categories using the same pairwise comparisons, filtering, and proxy/reference definitions as in the DA analysis, separately across model sizes and training checkpoints.

\begin{figure}[h]
    \centering
    \includegraphics[width=\linewidth]{figures/da_categories.png}
    \caption{
    \textbf{Evolution of pairwise intervention effects from proxy to reference.}
    We categorize each pairwise comparison according to whether the absolute performance gap between the two variants is within a tolerance of $\epsilon=0.5$ percentage points at the proxy and reference points.
    \emph{Tie--tie} indicates that the variants remain within the tolerance at both points; \emph{difference--tie} that a proxy-scale difference disappears at the reference; \emph{tie--difference} that a difference emerges only at the reference; and \emph{difference--difference} that the variants remain separated at both points.
    Left: category fractions across proxy model sizes relative to the 1.7B reference model.
    Right: category fractions across training checkpoints for the 350M model relative to its final checkpoint.
    }
    \label{fig:effect_categories}
\end{figure}


Figure~\ref{fig:effect_categories} shows that larger and later proxies become better at identifying whether an intervention remains consequential at the reference point. For DA-size, \emph{tie--difference} cases decrease with model size while \emph{difference--difference} cases increase, indicating that larger proxies are less likely to miss effects that persist at 1.7B. \emph{Difference--tie} remains comparatively stable, showing that some proxy-scale differences still wash out at the reference. The same pattern appears across training checkpoints, where \emph{tie--difference} decreases and \emph{difference--difference} increases with training progress. Thus, even when exact ranking recovery remains weak, larger models and later checkpoints provide a better signal of whether an intervention will matter at the reference point.